\paragraph{prime grammar、window-\texorpdfstring{$6$}{6} quotient 与 Lee--Yang 五重分支的三重不可压缩核}
在 Zeckendorf--G\"odel 稀疏素数语言的矩阵 Euler 刚性、window-\(6\) 的非自由群商障碍以及 Lee--Yang 分支塔的五个 \(T_2(E_{\min})\)-torsor 都已确定之后，还可把三者进一步压缩为一条统一的有限层阻碍链。其核心并不在于“无穷维极限尚未充分展开”，而在于下列三类有限层刚性已先行闭合：prime grammar 不能被标量 Euler 化，window-\(6\) 的 fold quotient 不能被自由轨道化，而 Lee--Yang 局部分支又不能被单链化为单一 dyadic sheet。

\begin{theorem}[window-\texorpdfstring{$6$}{6} 奇数残差的根系闭包障碍]\label{thm:xi-time-part9zp-window6-odd-residual-root-closure-obstruction}
\leanverified{paper\_xi\_time\_part9zp\_window6\_odd\_residual\_root\_closure\_obstruction}
沿用定理~\ref{thm:xi-window6-parity-charge-trigrading-b3} 的记号。记
\[
X_6=X_6^{\mathrm{cyc}}\sqcup X_6^{\mathrm{bdry}},
\qquad
|X_6|=21,\quad |X_6^{\mathrm{cyc}}|=18,\quad |X_6^{\mathrm{bdry}}|=3,
\]
并取表~\ref{tab:fold6_b3c3_root_dictionary_b3} 与表~\ref{tab:fold6_b3c3_root_dictionary_c3} 给出的 \(B_3/C_3\) 根字典，把 \(X_6^{\mathrm{cyc}}\) 与 \(18\) 个根一一对应。则：
\begin{enumerate}
  \item 全体 \(X_6\) 不可能识别为任何有限约化根系；
  \item 更强地，不存在任何有限约化根系 \(R\) 使
  \[
  X_6^{\mathrm{cyc}}\subseteq R,
  \qquad
  |R\setminus X_6^{\mathrm{cyc}}|=3.
  \]
\end{enumerate}
\end{theorem}
\begin{proof}
有限约化根系满足中心对称性
\[
\alpha\in R\Longrightarrow -\alpha\in R,
\]
且 \(0\notin R\)，故其基数必为偶数。

若 \(X_6\) 本身是某个有限约化根系，则其基数应为偶数；但
\[
|X_6|=21
\]
为奇数，矛盾。

再设存在有限约化根系 \(R\) 使
\[
X_6^{\mathrm{cyc}}\subseteq R,
\qquad
|R\setminus X_6^{\mathrm{cyc}}|=3.
\]
则
\[
|R|=|X_6^{\mathrm{cyc}}|+3=18+3=21,
\]
仍为奇数，同样矛盾。证毕。
\end{proof}

\begin{corollary}[window-\texorpdfstring{$6$}{6} 折叠商的非轨道化]\label{cor:xi-time-part9zp-window6-quotient-intrinsically-groupoidal}
设
\[
\omega\sim\eta
\iff
\Fold_6^{\mathrm{bin}}(\omega)=\Fold_6^{\mathrm{bin}}(\eta).
\]
则该等价关系既不能由任何有限群在 \(\Omega_6\) 上的自由作用轨道划分实现，也不能来自任何阿贝尔群的线性陪集划分。因而，window-\(6\) 的最小折叠商在有限层上已本质地属于群胚型压缩，而非自由群商。
\end{corollary}
\begin{proof}
定理~\ref{thm:xi-time-part67-window6-nonfree-group-quotient-obstruction} 已经证明：\(F_6=\Fold_6^{\mathrm{bin}}\) 的纤维分解既不能由自由群作用解释，也不能压平为阿贝尔线性陪集划分。故 \(\sim\) 的最小自然宿主只能允许稳定子随点变化，其群论语言必为群胚而非主群商。证毕。
\end{proof}

\begin{corollary}[Zeckendorf--G\"odel 素数语法的非标量 Euler 刚性]\label{cor:xi-time-part9zp-zg-prime-grammar-matrixization}
\leanverified{paper\_xi\_time\_part9zp\_zg\_prime\_grammar\_matrixization}
记
\[
\mathrm{ZG}
:=
\left\{
\prod_{i\in S}p_i:\ S\subset \NN\ \text{有限且不含相邻指标}
\right\},
\qquad
D_{\mathrm{ZG}}(s):=\sum_{n\in \mathrm{ZG}}n^{-s}.
\]
则 \(D_{\mathrm{ZG}}(s)\) 不能表示为任何纯局部标量 Euler 积
\[
\prod_{i\ge 1}F_i(p_i^{-s}),
\]
其中每个 \(F_i(z)\) 在 \(z=0\) 邻域解析且满足 \(F_i(0)=1\)。因此，prime-adjacency 约束迫使 \(\mathrm{ZG}\) 的 Euler 结构本质上是有序矩阵型，而不可能压缩为独立素点的标量局部乘积。
\end{corollary}
\begin{proof}
这正是定理~\ref{thm:xi-time-part62dgc-zg-matrix-euler-no-scalar-product} 的第三条。其证明表明：由于
\[
p_i\in \mathrm{ZG},
\qquad
p_{i+1}\in \mathrm{ZG},
\qquad
{p_i p_{i+1}}\notin \mathrm{ZG},
\]
对应指示函数 \(1_{\mathrm{ZG}}\) 不是乘法函数，故不存在保持同一 Dirichlet 系数的标量 Euler 积分解。证毕。
\end{proof}

\begin{theorem}[全支撑种子的面积--时间立方下界]\label{thm:xi-time-part9zp-allones-area-time-cubic-lower-bound}
\leanverified{paper\_xi\_time\_part9zp\_allones\_area\_time\_cubic\_lower\_bound}
设
\[
\Psi_m:\Omega_m\to X_m\times \Sigma^{\le L_m}
\]
为任意确定性可逆提升，字母表大小记为
\[
|\Sigma|=M\ge 2.
\]
再设
\[
\omega^{(m)}:=11\cdots 1\in\Omega_m
\]
是全 \(1\) 输入，并令 \(T(m)\) 为定理~\ref{thm:all-ones-reduction-length} 中 \(\omega^{(m)}\) 在规范化程序 \(\Pi_{m+1}\) 下的停机步数。则
\[
T(m)=\Bigl\lfloor \frac{m^2}{4}\Bigr\rfloor,
\]
并且
\[
\liminf_{m\to\infty}\frac{L_m\,T(m)}{m^3}
\ \ge\
\frac{1}{4}\cdot \frac{\log(2/\varphi)}{\log M}.
\]
\end{theorem}
\begin{proof}
由定理~\ref{thm:all-ones-reduction-length}，
\[
T(m)=\Bigl\lfloor \frac{m^2}{4}\Bigr\rfloor,
\qquad
\lim_{m\to\infty}\frac{T(m)}{m^2}=\frac14.
\]

另一方面，\(\Psi_m\) 为单射且满足
\[
\operatorname{pr}_{X_m}\circ \Psi_m=\Fold_m,
\]
故
\[
2^m=|\Omega_m|
\le
|X_m|\cdot |\Sigma^{\le L_m}|
=
F_{m+2}\sum_{\ell=0}^{L_m}M^\ell
<
F_{m+2}\,M^{L_m+1}.
\]
取 \(\log_M\) 得
\[
L_m\ge \left\lceil \log_M\!\Bigl(\frac{2^m}{F_{m+2}}\Bigr)\right\rceil-1,
\]
从而
\[
\liminf_{m\to\infty}\frac{L_m}{m}
\ge
\frac{\log(2/\varphi)}{\log M},
\]
因为 \(F_{m+2}=\Theta(\varphi^m)\)。

于是
\[
\liminf_{m\to\infty}\frac{L_m\,T(m)}{m^3}
\ge
\left(\liminf_{m\to\infty}\frac{L_m}{m}\right)
\left(\lim_{m\to\infty}\frac{T(m)}{m^2}\right)
\ge
\frac{\log(2/\varphi)}{\log M}\cdot \frac14.
\]
证毕。
\end{proof}

\begin{corollary}[Lee--Yang 分支塔的五重二维 dyadic 全移位分解]\label{cor:xi-time-part9zp-leyang-five-two-dimensional-fullshifts}
沿用定理~\ref{thm:xi-time-part62d-leyang-branch-inverse-limit-five-torsors} 的记号。记
\[
\Pi_n^{\mathrm{fin}}
=
\bigsqcup_{i=0}^{4}\Pi_{i,n},
\qquad
\Pi_{i,n}=P_i+E_{\min}[2^n].
\]
则：
\begin{enumerate}
  \item 对每个 \(i\) 与 \(n\)，\(\Pi_{i,n}\) 都与 \((\ZZ/(2^n\ZZ))^{2}\) 同构，因而 \(|\Pi_{i,n}|=4^n=2^{2n}\)；
  \item 过渡映射 \([2]:\Pi_{i,n+1}\to \Pi_{i,n}\) 在该识别下正是二坐标 dyadic 地址的“去最高位截断”；
  \item 分支逆极限 \(\mathcal R_\infty=\bigsqcup_{i=0}^{4}\mathcal R_\infty^{(i)}\) 非规范地分裂为五份 \(\ZZ_2^2\)；
  \item 若定义层熵
  \[
  h_{\mathrm{LY}}
  :=
  \lim_{n\to\infty}\frac1n\log \deg R(y_n),
  \]
  则
  \[
  h_{\mathrm{LY}}=\log 4=2\log 2.
  \]
\end{enumerate}
因此，Lee--Yang 局部分支动力学并非单一 dyadic 链，而是五份彼此平行的二维 dyadic 全移位。
\end{corollary}
\begin{proof}
由定理~\ref{thm:xi-time-part62d-leyang-branch-inverse-limit-five-torsors}，每个 \(\Pi_{i,n}\) 都是 \(E_{\min}[2^n]\) 的平移余类。又
\[
E_{\min}[2^n]\cong (\ZZ/(2^n\ZZ))^{2},
\]
故第 \(1\) 条成立。

推论~\ref{cor:xi-terminal-zm-leyang-finite-branch-regular-4ary-address} 已给出 \([2]\) 在有限分支层上的 regular \(4\)-ary 地址解释；在二坐标表示下，这正是两条独立 dyadic 坐标同时去最高位的截断映射，故得第 \(2\) 条。

第 \(3\) 条则是定理~\ref{thm:xi-time-part62d-leyang-branch-inverse-limit-five-torsors} 的非规范 torsor 识别
\[
\mathcal R_\infty^{(i)}\cong T_2(E_{\min})\cong \ZZ_2^2
\]
逐分支求不交并的结果。

最后，由
\[
\deg R(y_n)=5\cdot 4^n
\]
（见定理~\ref{thm:xi-time-part62d-leyang-branch-inverse-limit-five-torsors} 及定理~\ref{thm:xi-time-part62dgc-leyang-fivecore-dyadic-core-collapse}），直接计算得到
\[
h_{\mathrm{LY}}
=
\lim_{n\to\infty}\frac1n\log(5\cdot 4^n)
=
\log 4.
\]
证毕。
\end{proof}

\begin{conjecture}[grammar--quotient--sheet 三重不可压缩核]\label{conj:xi-time-part9zp-grammar-quotient-sheet-incompressible-kernel}
设某条沿 PRIMEGAME / Fold / Lee--Yang / moment-kernel 路线逼近 classical \(\zeta/\Xi\) 的有限化正规化，同时施行以下三类压缩：
\begin{enumerate}
  \item 将 prime grammar 标量化为独立 Euler 因子；
  \item 将 fold quotient 轨道化为自由群商或阿贝尔线性陪集商；
  \item 将 Lee--Yang 局部分支的五重 sheet 标号压缩为单分量 dyadic 链。
\end{enumerate}
则该正规化必然失败。其不可消去的根源分别在于：
\begin{enumerate}
  \item prime adjacency 迫使 Euler 结构本质上矩阵化；
  \item window-\(6\) 的折叠等价关系本质上不是自由轨道商，也不是线性陪集商；
  \item Lee--Yang 局部分支动力学本质上是五份平行的二维 dyadic 全移位，而非单链。
\end{enumerate}
若此猜想成立，则 classical \(\zeta/\Xi\) 路线中的真正困难，不首先表现为“尚未找到足够大的 ambient space”，而首先表现为：任何同时压缩 prime grammar、fold quotient 与局部分支标签的有限化范式，都会精确丢失三类彼此独立的刚性数据。
\end{conjecture}

\noindent
由此，PRIMEGAME/FRACTRAN 侧的非标量语法、window-\(6\) 侧的非自由 quotient，以及 Lee--Yang 侧的五重分支标签便不再只是三条互不相干的离散现象，而是同一条 obstruction chain 的三个有限层面：prime grammar 不能被单点 Euler 化，fold quotient 不能被主群商化，而局部分支又不能被单链化。若要继续沿 RH 路线推进，则更自然的方向已不再是寻找一个更大的单一算子去包揽全部数据，而是系统建立一条 grammar--quotient--sheet obstruction theory。

\endinput
